\section{Proof-Based Calculus \& Integration}

\subsection{The Mean Value Theorem and Rigorous Generalizations}

The transition from the topological foundations of $\mathbb{R}$ to differential calculus requires establishing the relationship between the local behavior of a function's derivative and its global variations over an interval. The cornerstone of this bridge is the Mean Value Theorem, which relies essentially on the topological compactness of closed intervals. 

Before proving the Mean Value Theorem, we first establish Rolle's Theorem, which itself depends on the Extreme Value Theorem (a direct consequence of the Heine-Borel theorem established in Chapter 1).

\begin{lemma}[Fermat's Theorem on Stationary Points]
Let $f: (a, b) \to \mathbb{R}$ be a function. If $f$ attains a local maximum or local minimum at a point $c \in (a, b)$, and if $f$ is differentiable at $c$, then $f'(c) = 0$.
\end{lemma}

\begin{proof}
Assume $f$ attains a local maximum at $c$. By definition, there exists an $\epsilon > 0$ such that for all $x \in (c - \epsilon, c + \epsilon)$, $f(x) \le f(c)$. 
For $x \in (c, c + \epsilon)$, $x - c > 0$ and $f(x) - f(c) \le 0$, yielding a difference quotient $\frac{f(x) - f(c)}{x - c} \le 0$. Taking the right-hand limit as $x \to c^+$, we find $f'(c) \le 0$.
Conversely, for $x \in (c - \epsilon, c)$, $x - c < 0$ and $f(x) - f(c) \le 0$, yielding $\frac{f(x) - f(c)}{x - c} \ge 0$. Taking the left-hand limit as $x \to c^-$, we find $f'(c) \ge 0$.
Since $f$ is differentiable at $c$, the left and right limits must coincide, forcing $f'(c) = 0$. The proof for a local minimum proceeds symmetrically.
\end{proof}

\begin{theorem}[Rolle's Theorem]
Let $f: [a, b] \to \mathbb{R}$ be continuous on the closed interval $[a, b]$ and differentiable on the open interval $(a, b)$. If $f(a) = f(b)$, then there exists at least one point $c \in (a, b)$ such that $f'(c) = 0$.
\end{theorem}

\begin{proof}
Because $[a, b]$ is compact and $f$ is continuous, the Extreme Value Theorem guarantees that $f$ attains its global maximum $M$ and global minimum $m$ on $[a, b]$. 
If $M = m$, then $f(x)$ is constant on $[a, b]$. Consequently, $f'(x) = 0$ for all $x \in (a, b)$, and any $c \in (a, b)$ satisfies the theorem.
If $M > m$, then since $f(a) = f(b)$, $f$ must attain at least one of its extrema at a point $c$ in the open interval $(a, b)$. By Fermat's Theorem, because $f$ is differentiable at this interior extremum, $f'(c) = 0$.
\end{proof}

\begin{theorem}[Lagrange Mean Value Theorem]
Let $f: [a, b] \to \mathbb{R}$ be continuous on $[a, b]$ and differentiable on $(a, b)$. There exists $c \in (a, b)$ such that:
\[
f'(c) = \frac{f(b) - f(a)}{b - a}
\]
\end{theorem}

\begin{proof}
We construct an auxiliary function $g: [a, b] \to \mathbb{R}$ representing the difference between $f(x)$ and the secant line passing through $(a, f(a))$ and $(b, f(b))$:
\[
g(x) = f(x) - f(a) - \frac{f(b) - f(a)}{b - a}(x - a)
\]
Since $f(x)$ and polynomials are continuous on $[a, b]$ and differentiable on $(a, b)$, $g(x)$ shares these properties. Evaluating $g$ at the endpoints yields:
\[
g(a) = f(a) - f(a) - 0 = 0
\]
\[
g(b) = f(b) - f(a) - (f(b) - f(a)) = 0
\]
Since $g(a) = g(b) = 0$, $g$ satisfies the hypotheses of Rolle's Theorem. Thus, there exists a point $c \in (a, b)$ such that $g'(c) = 0$. Differentiating $g(x)$ gives:
\[
g'(x) = f'(x) - \frac{f(b) - f(a)}{b - a}
\]
Setting $g'(c) = 0$ implies $f'(c) = \frac{f(b) - f(a)}{b - a}$, completing the proof.
\end{proof}

\begin{theorem}[Cauchy's Mean Value Theorem]
Let $f, g: [a, b] \to \mathbb{R}$ be continuous on $[a, b]$ and differentiable on $(a, b)$. There exists $c \in (a, b)$ such that:
\[
(f(b) - f(a))g'(c) = (g(b) - g(a))f'(c)
\]
\end{theorem}

\begin{proof}
Define the auxiliary function $h: [a, b] \to \mathbb{R}$ as follows:
\[
h(x) = (f(b) - f(a))g(x) - (g(b) - g(a))f(x)
\]
The function $h$ is continuous on $[a, b]$ and differentiable on $(a, b)$ due to the respective properties of $f$ and $g$. Evaluating $h$ at the endpoints:
\[
h(a) = f(b)g(a) - f(a)g(a) - g(b)f(a) + g(a)f(a) = f(b)g(a) - g(b)f(a)
\]
\[
h(b) = f(b)g(b) - f(a)g(b) - g(b)f(b) + g(a)f(b) = g(a)f(b) - f(a)g(b)
\]
Thus, $h(a) = h(b)$. By Rolle's Theorem, there exists $c \in (a, b)$ such that $h'(c) = 0$. Differentiating $h$ yields:
\[
h'(x) = (f(b) - f(a))g'(x) - (g(b) - g(a))f'(x)
\]
Setting $h'(c) = 0$ directly produces $(f(b) - f(a))g'(c) = (g(b) - g(a))f'(c)$.
\end{proof}

A powerful generalization of the Mean Value Theorem arises when analyzing higher-order derivatives, leading to Taylor's Theorem, which quantifies the error in polynomial approximations of functions.

\begin{theorem}[Taylor's Theorem with Lagrange Remainder]
Let $f: [a, b] \to \mathbb{R}$ be a function such that $f$ and its first $n$ derivatives are continuous on $[a, b]$, and $f^{(n+1)}$ exists on $(a, b)$. Let $x_0 \in [a, b]$. For any $x \in [a, b]$ with $x \neq x_0$, there exists a point $c$ strictly between $x$ and $x_0$ such that:
\[
f(x) = \sum_{k=0}^n \frac{f^{(k)}(x_0)}{k!}(x - x_0)^k + \frac{f^{(n+1)}(c)}{(n+1)!}(x - x_0)^{n+1}
\]
\end{theorem}

\begin{proof}
Fix $x \in [a, b]$. We define a real number $M$ such that the following equation holds:
\[
f(x) = \sum_{k=0}^n \frac{f^{(k)}(x_0)}{k!}(x - x_0)^k + M(x - x_0)^{n+1}
\]
Our goal is to show that $M = \frac{f^{(n+1)}(c)}{(n+1)!}$ for some $c$ between $x_0$ and $x$. To this end, define a function $F(t)$ on the closed interval bounded by $x_0$ and $x$:
\[
F(t) = f(x) - \sum_{k=0}^n \frac{f^{(k)}(t)}{k!}(x - t)^k - M(x - t)^{n+1}
\]
By construction, $F(x) = f(x) - f(x) - 0 = 0$, and by our definition of $M$, $F(x_0) = 0$. Since $F$ satisfies the conditions of Rolle's Theorem, there exists a point $c$ strictly between $x_0$ and $x$ such that $F'(c) = 0$. 

We compute the derivative of $F(t)$ with respect to $t$. The summation expands as a telescoping series when differentiated using the product rule:
\[
\frac{d}{dt} \left( \frac{f^{(k)}(t)}{k!}(x - t)^k \right) = \frac{f^{(k+1)}(t)}{k!}(x - t)^k - \frac{f^{(k)}(t)}{(k-1)!}(x - t)^{k-1}
\]
Summing this from $k=0$ to $n$, all terms cancel except the final one:
\[
\frac{d}{dt} \sum_{k=0}^n \frac{f^{(k)}(t)}{k!}(x - t)^k = \frac{f^{(n+1)}(t)}{n!}(x - t)^n
\]
Thus, the derivative of $F(t)$ simplifies to:
\[
F'(t) = - \frac{f^{(n+1)}(t)}{n!}(x - t)^n + M(n+1)(x - t)^n
\]
Evaluating at $t = c$ and setting $F'(c) = 0$ yields:
\[
0 = (x - c)^n \left( M(n+1) - \frac{f^{(n+1)}(c)}{n!} \right)
\]
Since $c \neq x$, $(x - c)^n \neq 0$, which mandates:
\[
M(n+1) = \frac{f^{(n+1)}(c)}{n!} \implies M = \frac{f^{(n+1)}(c)}{(n+1)!}
\]
Substituting this value of $M$ back into the original equation for $f(x)$ completes the proof.
\end{proof}
